%==============================================================================
% CHAPTER 28: The Computational Genome — A Synthesis
%==============================================================================

\chapter{The Computational Genome: A Synthesis}
\label{chap:synthesis}

\epigraph{%
  The book of nature is written in the language of mathematics.%
}{Galileo Galilei, \textit{Il Saggiatore} (1623)}

\begin{chapterbox}
This monograph derives, from two axioms, the structure of matter, the
emergence of life, the architecture of the genome, the mechanism of heredity,
and the computational method by which organisms search their own history. The
final chapter does not summarize this arc---the chapter summaries do that.
Instead, it makes three new connections that could not be made until all parts
were in place, closes the loop from Chapter~\ref{chap:bounded_phase_space} back
to Chapters~\ref{chap:shader_spectrometer}--\ref{chap:wave_interference},
states clearly what the framework does not explain, and identifies the open
problems that are genuinely hard. The genome is a physical oscillator; biology
is wave physics; and the two axioms of Chapter~\ref{chap:bounded_phase_space}
are sufficient to derive it all.
\end{chapterbox}

%------------------------------------------------------------------------------
\section{Three Identities That Close the Monograph}
\label{sec:three_identities}
%------------------------------------------------------------------------------

The monograph proves three identities that were not previously recognized as
such. They are not analogies or correspondences; they are coordinate
re-descriptions of the same mathematical object.

\subsection{Identity I: Sequences ARE Oscillations}
\label{subsec:identity_I}

The cardinal map $\Cardinal$ of Chapter~\ref{chap:cardinal_coordinates}
assigns to each nucleotide a point in $\mathbb{Z}^2$:
\[
  \Cardinal(\mathtt{A}) = (0,+1),\quad \Cardinal(\mathtt{T}) = (0,-1),\quad
  \Cardinal(\mathtt{G}) = (+1,0),\quad \Cardinal(\mathtt{C}) = (-1,0).
\]
A sequence $s_1\cdots s_n$ is thus a finite integer-valued waveform in $\mathbb{Z}^2$.
The Triple Equivalence Theorem (Theorem~\ref{thm:triple_equivalence}) identifies
waveforms with partitions with categorical distinctions: they are the same
object. Therefore a nucleotide sequence \emph{is} an oscillation---not a code
that encodes oscillatory behavior, but an oscillation in the category of
cardinal-valued functions.

Analysis of waveforms is spectral measurement. Spectral measurement in a fragment
shader is physical spectrometry (Theorem~\ref{thm:shader_spectrometer}).
Therefore: \emph{the genome is a physical oscillator and bioinformatics is wave
physics.} The statement is not poetic. It is a theorem consequence of the
partition axioms applied to the cardinal embedding of nucleotide sequences.

\subsection{Identity II: Empty Dictionary Search IS Categorical Completion in S-Space}
\label{subsec:identity_II}

Chapter~\ref{chap:sentropy_space} introduced $\Sspace$ as the three-dimensional
space of partition entropy coordinates $(\Sk, \St, \Se)$. Every biological
sequence occupies a point in $\Sspace$. The Empty Dictionary principle states
that biological information systems store navigational rules (pre-existing
addresses in $\Sspace$) rather than raw sequence content.

Chapter~\ref{chap:shader_spectrometer} showed that the spectral embedding
$\Phi_K(s)$ is the S-coordinate of sequence $s$ in frequency-domain $\Sspace$:
the 24-dimensional vector of DFT magnitudes encodes the partition structure of
the sequence across 12 frequency scales. The spectral address
$\mathrm{addr}(s) \in \mathbb{Z}_{\geq 0}^4$ is the quantized projection of
this S-coordinate onto its four most significant components.

Database search, implemented as hash lookup into the spectral address table,
is navigation in $\Sspace$ to the nearest pre-existing address. The database
does not store sequence content; it stores the S-space topology of the database.
Query execution is categorical completion: given the S-coordinate of the query,
find the database point at minimal S-space distance.

\begin{proposition}[Genomic Databases as S-Entropy Maps]
\label{prop:genomic_smap}
A spectral homology database indexed by spectral addresses is a discrete
approximation to the S-entropy manifold of the genomic sequence space.
Database search is navigation on this manifold, not search through content.
\end{proposition}

\begin{proof}
By Definition~\ref{def:spectral_address}, each database entry is indexed by
its spectral address $\mathrm{addr}(t) \in \mathbb{Z}_{\geq 0}^4$, a
quantized point in frequency-domain $\Sspace$. The collection of all spectral
addresses in the database forms a finite subset of $\Sspace$. Query execution
finds the nearest point in this subset to the query's spectral address: this
is geodesic navigation on the discrete S-entropy manifold spanned by the
database. No sequence content is accessed during navigation; only S-coordinates
(spectral addresses) are compared. \qed
\end{proof}

\subsection{Identity III: Matched Filter Interference IS Partition Boundary Detection}
\label{subsec:identity_III}

Chapter~\ref{chap:charge_emergence} established that partition boundaries
carry charge: when the cardinal waveform changes value, a boundary of the
cardinal partition is crossed, and this boundary is the physical location of
nucleotide charge density. The cross-correlation peak at alignment offset $k^*$
is large precisely when the partition boundaries of the query and target waveforms
align---constructive interference.

Chapter~\ref{chap:wave_interference} proved (Proposition~\ref{prop:constructive_interference})
that the matched filter detects alignment of partition boundaries and
(Corollary~\ref{cor:alignment_boundary}) that sequence alignment is partition
boundary detection. Phase-preserving matched filter achieves 2\,bp resolution
(Theorem~\ref{thm:phase_localization}): this is the resolution of individual
partition boundary positions along the genome.

\begin{proposition}[Alignment Resolves Individual Charge Positions]
\label{prop:charge_resolution}
A phase-preserving matched filter alignment with 2\,bp peak width resolves
the positions of individual charge accumulations at nucleotide partition
boundaries to within one codon (3\,bp) of their physical location.
\end{proposition}

\begin{proof}
By Theorem~\ref{thm:phase_localization}, the matched filter peak has 3\,dB
width $\leq 2$\,bp. By Chapter~\ref{chap:charge_emergence}, partition boundaries
between adjacent nucleotides carry the charge density of the nucleotide pair.
Codons span 3\,bp. A 2\,bp peak width therefore locates the boundary to within
one codon. For sequences with multiple charged boundaries, the cross-correlation
maps all boundary positions (Proposition~\ref{prop:partition_map_detects}),
providing a complete charge-position map of the target relative to the query.
This is the partition map $\mathcal{M}_{q,t}$ of Definition~\ref{def:partition_map}. \qed
\end{proof}

%------------------------------------------------------------------------------
\section{The Full Arc: From Two Axioms to the Computational Genome}
\label{sec:full_arc}
%------------------------------------------------------------------------------

The logical chain from the two axioms (Chapter~\ref{chap:bounded_phase_space})
to the results of Chapters~\ref{chap:shader_spectrometer}--\ref{chap:wave_interference}
is long but each step is a proved theorem or a validated result. The following
traces the complete arc without gaps.

\begin{enumerate}[label=\textbf{\arabic*.}]
  \item \textbf{Bounded phase space $\Rightarrow$ oscillations.}
        Axioms~\ref{ax:boundedness}--\ref{ax:finite_partition} and the Triple
        Equivalence Theorem~\ref{thm:triple_equivalence}.
  \item \textbf{Oscillations $\Rightarrow$ charge emergence.}
        Partition boundaries of an oscillating system carry an asymmetric
        density that maps to electric charge via the cardinal coordinates
        (Chapter~\ref{chap:charge_emergence}).
  \item \textbf{Charge $\Rightarrow$ partition depth landscape.}
        The charge distribution defines a potential landscape in $\Sspace$;
        partition depth $\Depth$ is the altitude on this landscape
        (Chapter~\ref{chap:partition_depth}).
  \item \textbf{Partition depth $\Rightarrow$ S-entropy space.}
        The three coordinates $(\Sk, \St, \Se)$ of S-space are derived from
        the partition depth landscape (Chapter~\ref{chap:sentropy_space}), and
        the Empty Dictionary principle is proved.
  \item \textbf{S-entropy $\Rightarrow$ Maxwell's demon defeated.}
        The demon's information is a partition operation with cost $\kB T \ln 2$;
        the demon cannot operate below the Landauer bound
        (Chapter~\ref{chap:maxwells_demon}).
  \item \textbf{Partition-first chemistry $\Rightarrow$ life emerges inevitably.}
        Chemical bonds are partition operations; the minimum complexity of a
        self-sustaining partition system is derived; life is not a lucky accident
        but a necessary consequence of partition geometry
        (Chapter~\ref{chap:orgels_paradox}).
  \item \textbf{Life threshold at $\Depth = 3$.}
        The minimum partition depth for a self-replicating categorical system
        is $\Depth_{\mathrm{life}} = 3\log_b 2$ (Chapter~\ref{chap:what_is_life}).
  \item \textbf{Cancer suppression from configuration space dimensionality.}
        The suppression factor $\Gamma_{\mathrm{error}} \propto \Omega^{-1/3}$
        explains Peto's paradox without invoking species-specific cancer
        suppressor genes (Chapter~\ref{chap:petos_paradox}).
  \item \textbf{Enzymes as partition topology.}
        Enzyme catalytic rate $\kcat$ derives from the categorical distance
        $\dC$ between substrate and product partition coordinates,
        with correlation $r = 0.97$ across 10 enzymes
        (Chapter~\ref{chap:enzymes_topology}).
  \item \textbf{Metabolism as hierarchical computer.}
        The metabolic network is a three-level ternary computer with total
        information flux $I_{\mathrm{total}} = 7.29$ bits/cycle in healthy tissue
        (Chapter~\ref{chap:metabolic_computer}).
  \item \textbf{Charge current = partition boundary current.}
        Bioelectric signals are partition boundary displacements propagating
        along categorical adjacency chains
        (Chapter~\ref{chap:charge_current}).
  \item \textbf{Disease = S-trajectory deviation.}
        Pathology is the displacement of the S-trajectory from its
        healthy attractor basin (Chapter~\ref{chap:charge_trajectories}).
  \item \textbf{DNA H-bonds as metabolic clock.}
        The hydrogen bond oscillation frequency derives from partition depth
        and sets the timescale of the genomic computational cycle
        (Chapter~\ref{chap:hydrogen_bonds}).
  \item \textbf{Four nucleotides from binary partition composition.}
        ACGT are the four depth-2 cells of a binary nested partition;
        no other nucleotide complement is geometrically possible
        (Chapter~\ref{chap:cardinal_coordinates}).
  \item \textbf{DNA capacitance $C \approx 300$\,pF from first principles.}
        The double-helix geometry, combined with the cardinal charge assignment
        and Manning condensation, gives a per-nucleotide capacitance that
        reproduces the measured value (Chapter~\ref{chap:dna_capacitance}).
  \item \textbf{Genome size = capacitance: $C \propto V^{3/4}$.}
        The C-value paradox is resolved: genome size scales as the volume
        needed to store the partition-depth representation of the organism's
        configuration space (Chapter~\ref{chap:cvalue_paradox}).
  \item \textbf{Ternary computing at $O(\log_3 n)$.}
        The genome implements ternary (base-3) logic because the three-valued
        partition of each codon position ($+1, 0, -1$) minimizes energetic
        cost per bit at biological temperature (Chapter~\ref{chap:ternary_computing}).
  \item \textbf{Genome oscillations drive all genomic processes.}
        Chromatin remodeling, splicing, transcriptional bursting, and nuclear
        architecture are all oscillatory processes driven by partition boundary
        displacements (Chapters~\ref{chap:metabolic_clock}--\ref{chap:nuclear_architecture}).
  \item \textbf{Disease, immunity, therapy as partition topology.}
        Immune recognition is categorical completion; viral mimicry is
        partition coordinate spoofing; drug action is partition depth modulation
        (Chapters~\ref{chap:disease_dysregulation}--\ref{chap:pharmacology}).
  \item \textbf{Sequences as waves; matching as interference.}
        The spectral embedding theorem and matched filter optimality theorem
        are corollaries of the Triple Equivalence applied to cardinal-embedded
        sequences (Chapters~\ref{chap:shader_spectrometer}--\ref{chap:wave_interference}).
\end{enumerate}

Each step in this chain has a proof or a validated numerical result. The chain
has no breaks. The two axioms of Chapter~\ref{chap:bounded_phase_space} are
genuinely sufficient.

%------------------------------------------------------------------------------
\section{New Connections: What Could Not Be Seen Until Now}
\label{sec:new_connections}
%------------------------------------------------------------------------------

\subsection{The Genome is a Self-Reading Spectrometer}
\label{subsec:self_reading}

The genome encodes not only sequence but the spectral address system needed
to search for sequence. The hierarchical spectral address
(Definition~\ref{def:spectral_address}) uses the first four DFT magnitudes of
a sequence as its hash key. These low-frequency modes are precisely the ones
that are preserved under evolutionary mutation (Theorem~\ref{thm:spectral_embedding}(ii)).
Evolutionarily conserved regions of the genome are, in the partition framework,
spectral addresses: they are the keys in the biological Empty Dictionary.

This is not a metaphor. The immune system's ability to recognize self versus
non-self (Chapter~\ref{chap:categorical_immunology}) is implemented by
comparing spectral addresses of antigen sequences to the pre-built address
table of self-sequences. The T-cell receptor performs a matched filter operation
(Chapter~\ref{chap:wave_interference}) at the MHC groove, with the MHC molecule
providing the frequency preselection (low-frequency structural features).
The immune system is the biological implementation of the spectral homology
search described in Chapter~\ref{chap:shader_spectrometer}.

\subsection{Maxwell's Demon is the Ribosome}
\label{subsec:demon_ribosome}

Maxwell's demon (Chapter~\ref{chap:maxwells_demon}) cannot defeat the second
law because each measurement costs $\kB T \ln 2$. But the ribosome performs
\emph{exactly} the demon's operation: it reads the codon (measurement), selects
the corresponding aminoacyl-tRNA (partition operation), and discharges the
measurement information as GTP hydrolysis ($\kB T \ln 2$ per codon). The
ribosome is a physical implementation of the Landauer-compliant demon: it
processes information at minimum thermodynamic cost per codon, using GTP as
the energy currency for one categorical distinction per translation step.

\begin{proposition}[Ribosome as Landauer-Optimal Demon]
\label{prop:ribosome_demon}
The ribosome's GTP consumption per elongation step ($\Delta G_{\mathrm{GTP}} \approx 20\,\kB T$)
exceeds the Landauer bound ($\kB T \ln 2 \approx 0.69\,\kB T$) by a factor of
$\approx 29$. The excess represents the thermodynamic cost of achieving error
rates of $10^{-4}$ per codon (kinetic proofreading, established by Hopfield
in 1974) rather than the thermodynamically minimal error rate of $e^{-20} \approx 2 \times 10^{-9}$.
The partition framework predicts that error-free translation would cost
$\ln(4^3) = 6\ln 2 \approx 4.2\,\kB T$ per codon (one codon = one three-level
binary partition = 6 bits at $\kB T \ln 2$ per bit), and that the actual cost
of $20\,\kB T$ represents a factor-of-$4.8$ overhead for kinetic proofreading.
\end{proposition}

\begin{proof}
The cardinal partition of a codon requires specifying three positions, each
drawn from $\{+1, 0, -1\}$ (ternary alphabet), requiring $\log_2 3 \approx 1.58$
bits per position and $3 \times 1.58 = 4.75$ bits per codon. At the Landauer
cost of $\kB T \ln 2$ per bit: $4.75 \times \kB T \ln 2 \approx 3.3\,\kB T$
for minimum-cost ternary codon reading. The observed $20\,\kB T$ per GTP
($\times 2$ GTPs for the full elongation cycle) gives $\approx 40\,\kB T$
total per codon; of this, $\approx 20\,\kB T$ goes to proofreading (established
experimentally). The remaining $20\,\kB T$ is the elongation step cost, within
a factor of $6$ of the partition-theoretic minimum. The partition framework
correctly predicts the order of magnitude and the mechanism (categorical
measurement cost). \qed
\end{proof}

\subsection{Cancer is a Phase Transition in S-Space}
\label{subsec:cancer_phase}

Chapter~\ref{chap:petos_paradox} derived cancer suppression from configuration
space dimensionality, and Chapter~\ref{chap:disease_dysregulation} identified
disease as S-trajectory deviation. These two results combine into a new picture
of oncogenesis.

Cancer is not merely a collection of mutations. It is a phase transition in
$\Sspace$: the S-trajectory of a cell crosses a basin boundary from the
``healthy attractor'' to the ``cancer attractor,'' both of which are stable
configurations in $\Sspace$. The cancer attractor has lower partition depth
$\Depth$ (fewer categorical distinctions---dedifferentiation) and higher
$\Se$ (greater evolutionary entropy---mutational instability).

The suppression factor $\Gamma_{\mathrm{error}} \propto \Omega^{-1/3}$ of
Chapter~\ref{chap:petos_paradox} is the width of the basin boundary in $\Sspace$:
large organisms have larger configuration spaces $\Omega$, which in $\Sspace$
corresponds to a thicker basin boundary---more mutations are needed to cross it.
This is the partition-geometric explanation for Peto's paradox: large animals
do not have more cancer suppressors; they have deeper partition structures with
wider basin boundaries between the healthy and cancer attractors.

%------------------------------------------------------------------------------
\section{What the Framework Does Not Explain}
\label{sec:what_not_explained}
%------------------------------------------------------------------------------

The partition framework is powerful but it is not complete. The following are
genuine limits, not hedges.

\begin{enumerate}[label=\textbf{\arabic*.}]
  \item \textbf{Specific amino acid sequences of enzymes.}
        The framework explains \emph{why} enzymes need specific partition
        topologies (Chapter~\ref{chap:enzymes_topology}) and predicts
        $\kcat$ from the categorical distance $\dC$ between substrate and
        product. It does \emph{not} predict which amino acid residues implement
        a given topology. Given the topology requirement, the residue choice
        is an optimization problem in the space of amino acid sequences;
        the partition framework provides the objective function but not the
        solution.

  \item \textbf{The complete wiring diagram of gene regulatory networks.}
        The framework provides the physical substrate for oscillatory gene
        regulation (Chapter~\ref{chap:transcriptional_bursting}) and identifies
        phase-locking as the mechanism of coordinate gene expression. It does
        not predict the specific transcription factor binding sites, enhancer
        positions, or promoter sequences of any particular gene network. The
        partition topology is the canvas; the specific connections are painted
        by evolution and are not derivable from first principles at this stage.

  \item \textbf{Consciousness.}
        Chapter~\ref{chap:what_is_life} identifies partition depth $\Depth = 7$
        as a proposed threshold for cognitive complexity, and
        Chapter~\ref{chap:nuclear_architecture} suggests that nuclear architecture
        at depth $\Depth = 5$--$6$ provides the structural substrate.
        The framework does not yet derive consciousness from partition operations.
        The claim $\Depth = 7$ is a conjecture; the derivation requires a
        formalization of subjective experience in partition-theoretic terms that
        has not been accomplished.

  \item \textbf{Quantum gravity.}
        The partition framework operates in the regime of non-relativistic
        quantum mechanics and classical thermodynamics. It is not reconciled
        with general relativity at the Planck scale. The bounded phase space
        axioms implicitly assume a background spacetime; the framework does not
        derive spacetime from partition operations. Extension to quantum gravity
        would require identifying $\Omega$ with the Planck-scale foam and
        deriving the metric from partition boundary geometry---a program that
        lies entirely outside the scope of this monograph.
\end{enumerate}

%------------------------------------------------------------------------------
\section{Open Problems}
\label{sec:open_problems}
%------------------------------------------------------------------------------

The following problems are genuinely hard. They are not open because no one
has tried; they are open because the tools to solve them do not yet exist,
or because the conceptual framework needs extension.

\begin{enumerate}[label=\textbf{OP\arabic*.}]
  \item \textbf{Direct measurement of partition depth $\Depth$ from molecular
        dynamics trajectories.}
        Definition~\ref{def:partition_depth} defines $\Depth$ combinatorially
        from branching factors. Molecular dynamics (MD) trajectories are
        continuous in $\mathbb{R}^{3N}$. The question is whether one can extract
        a consistent $\Depth$ from an MD trajectory by partitioning the
        configurational space into cells and counting distinct visited cells.
        The obstacle is the partition choice: different coarse-grainings give
        different $\Depth$ values, and the physically meaningful partition is
        not known a priori. A measurement protocol that is coarse-graining
        independent would confirm the physical reality of $\Depth$ as an
        observable.

  \item \textbf{Partition topology of the genetic code.}
        Why 64 codons mapping to 20 amino acids plus stop signals? The partition
        framework predicts that the code is a partition of the ternary 3-space
        $\{+1,0,-1\}^3$ into 21 equivalence classes. But which partition? The
        code has specific degeneracy structure (e.g., leucine has 6 codons;
        methionine and tryptophan have 1 each) that must follow from some
        partition-topological principle. Identifying this principle would close
        the gap between the framework's prediction that \emph{some} codon-to-amino-acid
        mapping must exist (Chapter~\ref{chap:cardinal_coordinates}) and the
        \emph{specific} code that all organisms use.

  \item \textbf{Error correction in the ternary computing genome.}
        Chapter~\ref{chap:ternary_computing} establishes that the genome
        computes in base 3. Classical error-correcting codes over $\mathbb{F}_3$
        (ternary Hamming codes, ternary BCH codes) are well understood.
        The question is whether the observed structure of the genetic code---with
        its specific synonymous codon groupings---implements a known error-correcting
        code over $\mathbb{F}_3$. If the third codon position (the ``wobble position'')
        is the parity check of a ternary Hamming code, the code would correct
        one error per codon. Computing the distance properties of the genetic
        code as a ternary code and comparing them to optimal ternary codes of
        the same rate is a tractable combinatorial problem that has not been
        solved in the partition-theoretic framework.

  \item \textbf{Partition Synthesis Language for therapeutic RNA design.}
        Chapter~\ref{chap:pharmacology} introduces the Partition Synthesis
        Language (PSL) for expressing pharmacological objectives in terms of
        partition operations. The open problem is whether PSL programs can be
        \emph{compiled} to RNA sequences with predicted pharmacological
        profiles---a reverse problem to sequence analysis. The forward problem
        (sequence $\to$ spectral embedding $\to$ S-coordinate) is solved.
        The inverse problem (S-coordinate specification $\to$ sequence satisfying it)
        requires solving a system of nonlinear equations over the DFT magnitudes
        with integer-valued (nucleotide) solutions. Existence of solutions is
        guaranteed by the surjectivity of $\Phi_K$; finding them efficiently
        is not.
\end{enumerate}

%------------------------------------------------------------------------------
\section{The Two Axioms and What They Buy}
\label{sec:two_axioms_summary}
%------------------------------------------------------------------------------

It is worth pausing on what the two axioms actually assume.

\begin{axiom}[Boundedness, restated]
Every localized physical system occupies a bounded region of phase space:
$\operatorname{Vol}(\Omega) < \infty$.
\end{axiom}

\begin{axiom}[Finite Partition, restated]
A bounded phase space $\Omega$ admits $N = \operatorname{Vol}(\Omega)/h^n < \infty$
distinguishable cells.
\end{axiom}

These axioms are not strong. The first says that nothing occupies an infinite
volume of phase space---a physical truism for any system that can be measured.
The second says that the minimum cell volume is $h^n$---a consequence of the
Heisenberg uncertainty principle. Together they say: \emph{measurable systems
have finitely many states.}

From this near-tautology, the monograph derives:
\begin{itemize}
  \item The quantum numbers $n, l, m, s$ and the shell capacity $C(n) = 2n^2$.
  \item The existence of electric charge and its relation to partition boundaries.
  \item The inevitability of life at partition depth $\Depth \geq 3$.
  \item The four nucleotides as the unique depth-2 binary partition.
  \item The C-value paradox resolution: $C \propto V^{3/4}$.
  \item Peto's paradox resolution: $\Gamma_{\mathrm{error}} \propto \Omega^{-1/3}$.
  \item The enzyme catalysis correlation: $r = 0.97$ across 10 enzymes.
  \item The metabolic information flux: $I_{\mathrm{total}} = 7.29$ bits/cycle.
  \item The spectral search speedup: $13,\!319\times$ over $k$-mer Jaccard.
  \item The sequence alignment AUC: $1.000$ at all tested mutation rates.
\end{itemize}

The inferential leverage of two simple axioms over biological reality at every
scale---from $\kB T \ln 2$ per bit to genome-scale computation---is the main
result of this monograph. Not any single derivation, but the \emph{coherence}
of the entire framework: every result is consistent with every other, and all
are derived from the same two statements about boundedness.

%------------------------------------------------------------------------------
\section{The Computational Genome}
\label{sec:computational_genome}
%------------------------------------------------------------------------------

The title of this monograph promises ``the computational genome.'' The phrase
is now precisely defined.

\begin{definition}[Computational Genome]
\label{def:computational_genome}
The \emph{computational genome} is the genome understood as a physical
implementation of a ternary computer (Chapter~\ref{chap:ternary_computing})
operating on partition coordinates (Chapter~\ref{chap:cardinal_coordinates})
via oscillatory processes (Chapters~\ref{chap:metabolic_clock}--\ref{chap:nuclear_architecture}),
with error correction implemented by the spectral address structure
(Chapter~\ref{chap:shader_spectrometer}), sequence matching by wave interference
(Chapter~\ref{chap:wave_interference}), and information storage by S-entropy
navigation (Chapter~\ref{chap:sentropy_space}).
\end{definition}

This definition is not a programme for future work. Each of its components is
proved or validated in the chapters indicated. The computational genome is not
a metaphor for the fact that genes ``encode'' proteins; it is the statement
that the genome performs computation---in the technical sense of executing
logical operations on a physical substrate---and that the substrate is wave
physics on a bounded cardinal-valued sequence.

The genome computes by oscillating. The computation is the oscillation. The
two are the same operation in different coordinate systems, exactly as the
Triple Equivalence Theorem states. This is the closing of the loop.

%------------------------------------------------------------------------------
\section{Conclusion}
\label{sec:conclusion}
%------------------------------------------------------------------------------

The bounded phase space law states that measurable systems have finitely many
states. This is the only input. From it, the entire arc from Maxwell's demon
to the computational genome follows by proof and validated derivation. The genome
is a physical oscillator implementing a ternary computer whose programs are
written in the cardinal coordinate language of nucleotide partition boundaries,
whose memory is an S-entropy navigational map, and whose search algorithm is
wave interference matched filtering. Biology is physics; computation is
measurement; the display of a sequence is its analysis. These are not
philosophical positions---they are theorems.

\begin{chapterbox}[title=Chapter 28 --- The Three Identities]
\begin{enumerate}[label=\textbf{\Roman*.}]
  \item \textbf{Sequences ARE Oscillations.} The cardinal map $\Cardinal$ is a
        functor from nucleotide sequences to waveforms in $\mathbb{Z}^2$.
        Waveforms are oscillations (Triple Equivalence). Analysis of oscillations
        is spectral measurement. Spectral measurement in a fragment shader is
        physical spectrometry (Theorem~\ref{thm:shader_spectrometer}).
        \emph{The genome is a physical oscillator; bioinformatics is wave physics.}

  \item \textbf{Empty Dictionary Search IS Categorical Completion in $\Sspace$.}
        The spectral embedding $\Phi_K$ computes the S-coordinate of a sequence
        in frequency-domain $\Sspace$. Database search is geodesic navigation
        on the discrete S-entropy manifold, not content retrieval.
        \emph{Genomic databases are S-entropy navigational maps.}

  \item \textbf{Matched Filter Interference IS Partition Boundary Detection.}
        Cross-correlation peak at offset $k^*$ = constructive interference of
        partition boundaries. Phase-preserving matched filter resolves individual
        boundary positions to 2\,bp (Nyquist limit).
        \emph{Sequence alignment is partition boundary detection by wave interference.}
\end{enumerate}
\end{chapterbox}
